\documentclass[10pt,dvipdfmx,m,compress]{beamer}

%\hypersetup{pdfpagemode=FullScreen}
\beamertemplatenavigationsymbolsempty
\usetheme{NIT}
\DeclareGraphicsExtensions{.jpg,.eps,.png,.pdf,.mps,.gif}

\usepackage{pgf,pgfarrows,pgfnodes,pgfautomata,pgfheaps}
\usepackage{amsmath}
\usepackage{amsfonts}
\usepackage[dvipdfmx]{graphicx}
\usepackage{color}
\usepackage{listings}
\lstset{basicstyle=\ttfamily\footnotesize}

%%%%%%%%%%%%%%%%%%%%%%%%%%%%%%%%%%%%%%%%%%%%%%%%%%%%%%%%%%%%%%%%%%%%%%%%%%%%%%%%%%%%%%%
% color definitions %%%%%%%%%%%%%%%%%%%%%%%%%%%%%%%%%%%%%%%%%%%%%%%%%%%%%%%%%%%%%%%%%%%
%%%%%%%%%%%%%%%%%%%%%%%%%%%%%%%%%%%%%%%%%%%%%%%%%%%%%%%%%%%%%%%%%%%%%%%%%%%%%%%%%%%%%%%
\definecolor{NITMainColor}{RGB}{255,130,0}
\definecolor{NITReverseColor}{RGB}{240,240,240}
\definecolor{NITgray}{RGB}{190,190,190}

%%%%%%%%%%%%%%%%%%%%%%%%%%%%%%%%%%%%%%%%%%%%%%%%%%%%%%%%%%%%%%%%%%%%%%%%%%%%%%%%%%%%%%%
% title page definition %%%%%%%%%%%%%%%%%%%%%%%%%%%%%%%%%%%%%%%%%%%%%%%%%%%%%%%%%%%%%%%
%%%%%%%%%%%%%%%%%%%%%%%%%%%%%%%%%%%%%%%%%%%%%%%%%%%%%%%%%%%%%%%%%%%%%%%%%%%%%%%%%%%%%%%
\setbeamercovered{transparent}
\setbeamerfont{author}{family=\rmfamily}

\author[Ryo KOBAYASHI]{Ryo KOBAYASHI}
\title[NIT Theme]{NIT Theme}
\institute{Department of Scientific and Engineering Simulation,\\
Nagoya Insititute of Technology}
\date[\today]{\today}
\titlegraphic{\includegraphics[height=1.5cm]{NIT_Logo.eps}}

% table of contents depth
\setcounter{tocdepth}{2}

\begin{document}
\begin{frame}[plain,b]
 \titlepage
\end{frame}

\usebackgroundtemplate{\includegraphics[width=\paperwidth]{NIT_bg.eps}}
\frame{\frametitle{Outline}\tableofcontents}

 \section{Installation}
  \subsection{How to use this theme?}
  \begin{frame}{What is the NIT?}
   \begin{columns}
    \begin{column}{.4\textwidth}
     \begin{figure}[h]
      \begin{center}
       \includegraphics[keepaspectratio,width=.8\textwidth]{NIT_logo.eps}
       \caption{The logo of NIT.}
      \end{center}
     \end{figure}
    \end{column}
    %
    \begin{column}{.6\textwidth}
     \begin{itemize}
      \item NIT is the acronym of \structure{\itshape Nagoya Institute of Technology}.
      \item The author is \structure{\itshape Ryo KOBAYASHI}, who is an assistant professor at NIT.
      \item Please contact to the author by e-mailing to \structure{\itshape ryo.kbys[at]gmail.com}.
            \vspace{1zw}
      \item This theme is not {\itshape Nagoya} (the city name where NIT exists),
            because there might be somebody who wants to make own theme in other institute in Nagoya.
     \end{itemize}
    \end{column}
   \end{columns}
  \end{frame}

  \begin{frame}{How to use this theme?}
   \begin{enumerate}
    \item Put beamer*NIT.sty and NIT*eps files to the place where the environment TEXINPUTS looks for.
    \item Call this theme by puting \listinline{\textbackslash{}usetheme\{NIT\}} in the preamble.
    \item Copy preambles of this source file to your own source file, especially color definitions are important for how it looks.
    \item Add ``compress'' to the option of beamer documentclass to keep the headline height fixed even if there are several subsections.
    \item After the title frame, put usebackgroundtemplate and specify {\itshape NIT\_bg.eps} as a background image.
   \end{enumerate}
  \end{frame}

 \section{Customization}
  \subsection{How to change logo images?} % subtitles are not used in this layout but still required for the header outline bullets
  \begin{frame}{How to change background and logo images?}
   \begin{itemize}
    \item The background and logo images are NIT\_bg.eps and NIT\_logo.eps, respectively.
    \item Modifying those EPS images results directly in how it looks.
    \item Title bar in each frame is included in the NIT\_bg.eps. Thus if you want to change frametitle, modify the NIT\_bg.eps file.
   \end{itemize}
  \end{frame}
  
  \subsection{How to customize colors?}
  \begin{frame}{How to customize the main color?}
   \begin{itemize}
    \item Change color definition at the preamble, {\itshape NITMainColor}.
    \item Modify the background image, {\itshape NIT\_bg.eps}, to change the color of the bottom of the frametitle banner.
   \end{itemize}
  \end{frame}

\end{document}
